Write $\C$ for $\Cl(0,1) = \mathrm{span}\{1, e_{1}\}$, a commutative field (\claimref{std-algebra.complex-numbers}), with conjugation $\overline{a + be_{1}} = a - be_{1}$, a ring automorphism fixing $\R$, and $z\bar{z} = a^{2} + b^{2}$.

\emph{(i), step 1: odd degree.} Let $p = X^{n} + \sum_{i < n}c_{i}X^{i}$, $n$ odd, and $M = 1 + \sum_{i < n}|c_{i}|$. For $|x| \ge M \ge 1$, $|\sum_{i<n}c_{i}x^{i}| \le (M - 1)|x|^{n-1} < |x|^{n}$, so $p(M) > 0 > p(-M)$. On $[-M, M]$, $|p(x) - p(y)| \le \Lambda|x - y|$ with $\Lambda = \sum_{l}l\,|c_{l}|M^{l-1}$ ($c_{n} = 1$), since $x^{i} - y^{i} = (x - y)\sum_{j}x^{j}y^{i-1-j}$. Let $s = \sup\{x \in [-M, M] : p(x) < 0\}$, which exists by completeness. If $p(s) < 0$, then $s < M$ and $p < 0$ on $[s, s + |p(s)|/\Lambda)$, against $s$ being an upper bound; if $p(s) > 0$, then $p > 0$ on $(s - p(s)/\Lambda, s]$, so $s - p(s)/\Lambda$ is a smaller upper bound. Hence $p(s) = 0$. The same argument applied to $X^{2} - a$, $a > 0$, on $[0, M]$ with $M = 1 + a$ (where $p(0) = -a < 0 < p(M)$) gives $\sqrt{a} \in \R$.

\emph{Step 2: square roots in $\C$.} For $z = a + be_{1}$ let $r = \sqrt{a^{2} + b^{2}} \ge |a|$, $x = \sqrt{(r + a)/2}$, $y = \sigma\sqrt{(r - a)/2}$ with $\sigma = 1$ if $b \ge 0$ and $-1$ otherwise. Then $(x + ye_{1})^{2} = x^{2} - y^{2} + 2xye_{1} = a + \sigma\sqrt{r^{2} - a^{2}}\,e_{1} = a + be_{1}$. So every monic quadratic $X^{2} - sX + P$ over $\C$ factors as $(X - \tfrac12(s + w))(X - \tfrac12(s - w))$ with $w^{2} = s^{2} - 4P$.

\emph{Step 3: Laplace's induction.} Write $n = \deg p = 2^{m}u$ with $u$ odd; we show by induction on $m$ that $p$ has a root in $\C$. For $m = 0$ this is step 1. For $m \ge 1$, let $L$ be a splitting field of $p$ over the field $\C = \Cl(0,1)$ (obtained by adjoining roots of irreducible factors one at a time, $K[X]/(g)$ being a field for $g$ irreducible over a field $K$), so that $\R \subseteq \C \subseteq L$. Let $r_{1}, \dots, r_{n} \in L$ be the roots of $p$ with multiplicity. For $t \in \R$ put $g_{t}(X) = \prod_{i<j}(X - (r_{i} + r_{j} + t\,r_{i}r_{j}))$, of degree $n(n-1)/2 = 2^{m-1}u(n - 1)$ with $u(n - 1)$ odd. Its coefficients are polynomials over $\R$ in the $r_{i}$, invariant under permutations of the $r_{i}$, hence (fundamental theorem on symmetric polynomials: by induction on the leading monomial in lexicographic order, such a polynomial is a polynomial in the elementary symmetric ones) polynomials over $\R$ in the coefficients of $p$: $g_{t}$ has real coefficients. By induction it has a root in $\C$, so for each $t$ some pair $i < j$ has $r_{i} + r_{j} + t\,r_{i}r_{j} \in \C$. There are infinitely many $t$ and finitely many pairs, so one pair serves two values $t \ne t'$; subtracting, $(t - t')r_{i}r_{j} \in \C$, so $P = r_{i}r_{j} \in \C$ and $s = r_{i} + r_{j} \in \C$. Then $r_{i}$ is a root of $X^{2} - sX + P$, whose roots in the field $L$ are those of step 2, which lie in $\C$. Finally, if $q$ is irreducible over $\R$ of degree at least $2$, it has a root $z \in \C \setminus \R$, also $\bar{z}$ is a root, and $(X - z)(X - \bar{z}) = X^{2} - (z + \bar{z})X + z\bar{z}$ has real coefficients and divides $q$ over $\C$, hence over $\R$ (division with remainder over $\R$ gives the same quotient and remainder as over $\C$); so $\deg q = 2$. The check verifies the square-root formula on squares of rational elements of $\Cl(0,1)$, the sign change of three odd-degree polynomials at $\pm M$, and for $p = (X^{2} + 1)(X^{2} + 4)$, with roots $\pm e_{1}$, $\pm 2e_{1}$ in $\Cl(0,1)$, that $g_{t}$ has real coefficients for five values of $t$ and equals $X^{6} + 10X^{4} + 9X^{2}$ at $t = 0$.

\emph{(ii)} A nonzero $x \in D$ is invertible: $y \mapsto xy$ is injective and linear on the finite-dimensional $D$, hence onto, so $xy = 1$ for some $y$, and $x(yx - 1) = 0$ gives $yx = 1$. For $x \in D$ the powers $1, x, \dots, x^{\dim D}$ are dependent, so $x$ has a monic minimal polynomial $m_{x}$ over $\R$; it is irreducible, since $m_{x} = gh$ with smaller degrees would give $g(x)h(x) = 0$ with both factors nonzero. By (i) its degree is $1$ or $2$. If $x \notin \R 1$, $m_{x} = X^{2} - 2bX + c$ with $b^{2} - c < 0$ (no real root), and $(x - b)^{2} = b^{2} - c < 0$, so $x - b \in V$: $D = \R 1 + V$. A real number with non-positive square is $0$, so $\R 1 \cap V = 0$. For $v \in V \setminus 0$, $v^{2} < 0$ (as $v^{2} = 0$ would make $v$ a zero divisor) and $v \notin \R 1$. $V$ is closed under real multiples, $(\lambda v)^{2} = \lambda^{2}v^{2}$. For $u, v \in V$: if $u = 0$ there is nothing to show; if $1, u, v$ are dependent and $u \ne 0$, then $v = \alpha + \beta u$ and $v^{2} = \alpha^{2} + \beta^{2}u^{2} + 2\alpha\beta u$ is real, so $\alpha\beta = 0$; $\beta = 0$ gives $v = \alpha \in \R 1 \cap V = 0$, and otherwise $u + v = (1 + \beta)u \in V$. If $1, u, v$ are independent, $u \pm v \notin \R 1$, so $(u + v)^{2} = \alpha(u + v) + \beta$ and $(u - v)^{2} = \gamma(u - v) + \delta$; adding, $2u^{2} + 2v^{2} = (\alpha + \gamma)u + (\alpha - \gamma)v + \beta + \delta$, and independence gives $\alpha = \gamma = 0$, so $(u + v)^{2} = \beta \in \R$. If $\beta > 0$, $(u + v - \sqrt{\beta})(u + v + \sqrt{\beta}) = 0$ would put $u + v \in \R 1$; $\beta = 0$ would give $u + v = 0$; so $\beta < 0$ and $u + v \in V$. Thus $V$ is a subspace and $D = \R 1 \oplus V$. For $u, v \in V$, $uv + vu = (u + v)^{2} - u^{2} - v^{2} \in \R$, so $B(u, v) = \tfrac12(uv + vu)$ is a symmetric bilinear form with $B(v, v) = v^{2} = Q(v) < 0$ for $v \ne 0$. Rejection $x \mapsto x - B(x, u)u^{-1}$ and division by $\sqrt{-Q(u)}$ (step 1) give an orthonormal basis $e_{1}, \dots, e_{k}$ of $V$: $e_{i}^{2} = -1$, $e_{i}e_{j} = -e_{j}e_{i}$ for $i \ne j$ (Proposition \claimref{std-algebra.orthogonal-anticommute}). By the universal property the inclusion of $V$ extends to a homomorphism $\phi: \Cl(0,k) \to D$, onto since its image contains $1$ and $V$; $\dim \Cl(0,k) = 2^{k}$ (Theorem \claimref{std-algebra.blade-basis}).

Let $k \ge 2$ and $u = e_{1}e_{2} \in D$. Then $u^{2} = -e_{1}^{2}e_{2}^{2} = -1$, so $u \in V$, and $ue_{1} + e_{1}u = e_{1}e_{2}e_{1} + e_{1}e_{1}e_{2} = e_{2} - e_{2} = 0$, likewise $ue_{2} + e_{2}u = 0$: $B(u, e_{1}) = B(u, e_{2}) = 0$. Writing $u = \sum_{i}c_{i}e_{i}$, $B(u, e_{j}) = -c_{j}$. If $k = 2$ this gives $u = 0$, against $u^{2} = -1$. If $k \ge 3$ and $j \ge 3$, $e_{j}$ anticommutes with $e_{1}$ and $e_{2}$, hence commutes with $u$, and $2ue_{j} = ue_{j} + e_{j}u = 2B(u, e_{j}) = -2c_{j}$ is real, so $e_{j} = u^{-1}(ue_{j}) = c_{j}u$ (as $u^{-1} = -u$). The vectors $e_{3}, \dots, e_{k}$ are independent and all lie in $\R u$, so $k \le 3$, and for $k = 3$, $-1 = e_{3}^{2} = -c_{3}^{2}$ gives $e_{3} = \pm e_{1}e_{2}$. Hence: $k = 0$ and $D = \R$; or $k = 1$ and $\phi: \Cl(0,1) \to D$ is onto between spaces of dimension $2$, an isomorphism; or $k = 3$, $D = \mathrm{span}\{1, e_{1}, e_{2}, e_{1}e_{2}\}$ is the image of the subalgebra generated by $e_{1}, e_{2}$, a copy of $\Cl(0,2)$ of dimension $4$, so $D \cong \Cl(0,2) \cong \HH$ (\claimref{std-algebra.quaternions}). In $\Cl(0,3)$ itself the element $ue_{3} = e_{1}e_{2}e_{3}$ is central with square $+1$, and $(1 - ue_{3})(1 + ue_{3}) = 0$: the map to $D$ sends it to $\pm 1$, the quaternionic quotient. That $\R$, $\Cl(0,1)$, $\Cl(0,2)$ have no zero divisors is Theorem \claimref{transport.exceptional-division-algebras}. The check verifies in $\Cl(0,1)$ and $\Cl(0,2)$ that an element squares to a real exactly when its real or its pure part vanishes and that $x^{2} - 2ax + (a^{2} - v^{2}) = 0$; in $\Cl(0,k)$, $k = 2, 3, 4$, the relations of $u = e_{1}e_{2}$ with the generators and $(1 - ue_{j})(1 + ue_{j}) = 0$; that $e_{1}, e_{2}, e_{1}e_{2}$ satisfy the relations of three generators of $\Cl(0,3)$; that $e_{1}e_{2} \notin \mathrm{span}\{1, e_{1}, e_{2}\}$; and the dimensions $1, 2, 4$.

\emph{(iii)} By Proposition \claimref{std-algebra.witt-basis-car} the $v_{i}$, $w_{i}$ satisfy $v_{i}^{2} = w_{i}^{2} = 0$, $v_{i}w_{i} + w_{i}v_{i} = 1$, and anticommute across different indices, and the $v_{i}w_{i}$ are commuting idempotents, so $f$ is idempotent. Each $v_{i}w_{i}$ commutes with $v_{j}$, $w_{j}$ for $j \ne i$ (two sign changes), so $v_{i}f = 0$ and $fw_{i} = 0$ (move the factor $v_{i}w_{i}$ next to $v_{i}$, respectively $w_{i}$, and use $v_{i}^{2} = w_{i}^{2} = 0$). Consider $fv_{J}^{\leftarrow}w_{K}f$. If $i \in K \setminus J$, move $w_{i}$ to the left through the $v_{j}$, $j \in J$, and the $w$'s before it, all anticommuting with it, to meet $f$: the product is $0$. If $i \in J \setminus K$, move $v_{i}$ to the right likewise: $0$. If $J = K$, the innermost pair $v_{j}w_{j}$ is even and commutes with the remaining letters, so $v_{J}^{\leftarrow}w_{J} = \prod_{j \in J}v_{j}w_{j}$, and $\prod_{j \in J}(v_{j}w_{j})f = f$. So $fv_{J}^{\leftarrow}w_{K}f = \delta_{JK}f$ and $E_{I,J}E_{K,L} = \delta_{JK}E_{I,L}$. By the same moves $E_{I,I} = \prod_{i \in I}w_{i}v_{i}\prod_{i \notin I}v_{i}w_{i}$, and $\sum_{I}E_{I,I} = \prod_{i}(w_{i}v_{i} + v_{i}w_{i}) = 1$. So the $E_{I,J}$ are $2^{k} \times 2^{k}$ matrix units; by Proposition \claimref{std-algebra.matrix-units} $\Cl(k,k) \cong \Mat_{2^{k}}(f\Cl(k,k)f)$, and comparing dimensions ($4^{k}$, Theorem \claimref{std-algebra.fock-representation}) the corner $f\Cl(k,k)f$ is $\R f$. Now $a \mapsto L(a)$ is an injective unital homomorphism $A \to \Mat_{d}(\R)$ ($L(a)$ applied to the coordinates of $1$ gives those of $a$), and $M \mapsto \sum_{s,t}M_{st}E_{I_{s},I_{t}}$ is an injective homomorphism $\Mat_{d}(\R) \to \Cl(k,k)$ (the $E$'s multiply as matrix units and are linearly independent) sending the identity to $\epsilon$; the composite is $\varphi$, and $\varphi(A) = \epsilon\varphi(A)\epsilon$. The check verifies the matrix-unit relations for $k = 1, 2$, $\sum_{I}E_{I,I} = 1$ for $k = 1, 2, 3$, and the embedding (homomorphism on basis pairs, injectivity, the idempotent unit) for the dual numbers, the split-complex numbers, $\R[x]/(x^{3})$, the upper triangular $2 \times 2$ matrices and $\Cl(0,2)$.

\emph{(iv)} $\epsilon_{\pm}^{2} = \tfrac14(1 \pm 2e_{1} + e_{1}^{2}) = \epsilon_{\pm}$, $\epsilon_{+}\epsilon_{-} = \tfrac14(1 - e_{1}^{2}) = 0$, $\epsilon_{+} + \epsilon_{-} = 1$, so $\Cl(1,0) = \R\epsilon_{+} \oplus \R\epsilon_{-} \cong \R \oplus \R$. In $\Cl(1,1)$, $n = e_{1} + e_{2}$ has $n^{2} = 1 - 1 = 0$, so $\mathrm{span}\{1, n\}$ is closed, $(a + bn)(c + dn) = ac + (ad + bc)n$, and $\rho \mapsto n$ is an isomorphism from $\Cl(0,0,1)$. The complex numbers and quaternions are Propositions \claimref{std-algebra.complex-numbers} and \claimref{std-algebra.quaternions}, and $\Cl(k,k) \cong \Mat_{2^{k}}(\R)$ is (iii) with Theorem \claimref{std-algebra.fock-representation}. The final sentence restates (ii) and (iii). The check verifies the idempotents of $\Cl(1,0)$, the dual-number product on a null vector of $\Cl(1,1)$, and $x\bar{x} = $ sum of squares in $\Cl(0,1)$ and $\Cl(0,2)$.
